% Copy this entire project in Overleaf; edit this file for a new document.
\documentclass[11pt,letterpaper,twoside]{article}
\usepackage[margin=1in,headheight=40pt]{geometry}

% NOTATION: choose one line. Add prooftrees only when using bussproofs trees.
\usepackage{andyspackage}
% \usepackage[kripke]{andyspackage}
% \usepackage[kripke,prooftrees]{andyspackage}

% STYLE: leave exactly one of these lines uncommented.
\usepackage[header=book,font=times,accent=purple]{andysnotes}
% \usepackage[header=plain,font=times,accent=purple]{andysnotes}
% \usepackage[header=none,font=times,accent=purple]{andysnotes}
% \usepackage[header=book,font=latinmodern,accent=black,sectionrules=false]{andysnotes}
% \usepackage[header=book,font=times,bibliography=angus]{andysnotes}
% Add showsectionnumbers=true to print main section numbers.
% If bibliography=angus is active, add after the selected style line:
% \addbibresource{docs/references.bib}

% GRAPHICS: uncomment only when you need figures, diagrams, or plots.
% \usepackage{andyspics}

% The title needs only these two fields.
\title{A Working Title}
\author{Andrew M. Hagy}
\date{}
% Add any fields that suit this project:
% \AndySubtitle{A subtitle}
% \AndyCourse{PHIL 4770}
% \AndyAssignment{Homework 3}
% \AndyDue{Friday, October 2, 2026}
% \AndyAffiliation{University of Pennsylvania}
% \AndyRunningTitle{Short title for even-page headers}

\begin{document}
\setcounter{page}{0} % change to 1 if desired

\AndyMakeTitle
% \AndyMakeTitlePage % use this INSTEAD of \AndyMakeTitle for a separate page
% \pagenumbering{roman} % before contents if front matter uses Roman numerals
% \tableofcontents
% \AndyStartBodyPages % after front matter: new page, Arabic page 0

\section{First section}
Main section numbers are hidden, but the counter still runs. In particular,
$\Gamma\entails\varphi$ means semantic consequence, while
$\Gamma\proves\varphi$ means derivability.

\subsection{A first subsection}
This is subsection 1.1. You can write $\RR$ or $\Reals$ for the reals.

\begin{DefinitionBlock}[Sample]
A set $A$ is a subset of $B$ when $A\containedin B$.
\end{DefinitionBlock}

\begin{aside}
An unboxed aside follows the same statement sequence.
\end{aside}

\begin{TheoremBlock}[Sample theorem]
If $A\containedin B$ and $B\containedin C$, then $A\containedin C$.
\begin{ProofBlock}
Take $x\in A$. Then $x\in B$ and $x\in C$.
\end{ProofBlock}
\end{TheoremBlock}

\begin{CorollaryBlock}
In particular, $A\containedin A$.
\end{CorollaryBlock}

% For a short reference list without a separate file:
% \begin{thebibliography}{9}
% \bibitem{key} Author, \emph{Title}, publication details.
% \end{thebibliography}
% Or, with bibliography=angus: \printbibliography
\end{document}
